\documentclass[11pt]{article}

\usepackage[utf8]{inputenc}
\usepackage[T1]{fontenc}
\usepackage{amsmath, amssymb}
\usepackage{graphicx}
\usepackage{booktabs}
\usepackage{hyperref}
\usepackage{geometry}
\usepackage{caption}
\usepackage{xcolor}
\usepackage{algorithm}
\usepackage{algorithmic}

\geometry{margin=2.5cm}

\title{Frustrated Synchronization in Doubt-Modulated FitzHugh--Nagumo Networks:\\
Attractor Diversity and Topological Phase Transitions}

\author{
    \textbf{Julien Chauvin} \\
    \small{\texttt{contact@cafevirtuel.org}} \\
    \small{Repository: \url{https://github.com/cafe-virtuel/mem4ristor-v2}}
}

\date{\today}

\begin{document}

\maketitle

% ============================================================================
% ABSTRACT (~150 words)
% ============================================================================
\begin{abstract}
Coupled oscillator networks generically synchronize, collapsing attractor diversity. We present a modified FitzHugh--Nagumo model in which a third state variable, \emph{doubt} ($u$), modulates coupling polarity via a shifted-tanh kernel $f(u) = \tanh(\pi(0.5-u)) + \delta$. Combined with a fraction $\eta$ of nodes receiving inverted stimulus (heretic nodes) and Gaussian noise for symmetry breaking, this mechanism sustains multi-state attractors from homogeneous initial conditions. On regular lattices ($N = 100$), the sustained attractor entropy stabilizes at $H_{\text{stable}} \approx 0.92 \pm 0.04$ bits (Shannon entropy over 5 discrete states, last 25\% of $3000$-step runs, $7$ seeds). Degree-normalized coupling ($D/\deg(i)^\gamma$) recovers diversity on sparse scale-free networks (Barab\'asi--Albert $m \leq 3$), but a systematic sweep over $m \in [1, 15]$ and $\gamma \in [0, 1]$ reveals an abrupt phase transition at $m \approx 5$: for denser topologies, no degree-based normalization preserves diversity ($H < 0.03$). The algebraic connectivity $\lambda_2$ predicts regime membership, suggesting that local coupling corrections fail when path redundancy exceeds a critical threshold.
\end{abstract}

% ============================================================================
% 1. INTRODUCTION
% ============================================================================
\section{Introduction}

Synchronization is a generic outcome of coupled oscillator dynamics~\cite{pikovsky2001,strogatz2000}. The Kuramoto model~\cite{acebrn2005}, distributed consensus algorithms~\cite{olfati2007}, and voter models~\cite{castellano2009} all converge toward uniform states, losing the diversity of initial conditions. While synchronization is desirable in many applications, it becomes a failure mode when heterogeneous steady states carry functional value---as in neural coding, opinion dynamics, or collective decision-making.

The problem is well characterized: on networks with homogeneous coupling, the largest Lyapunov exponent of the synchronized manifold is typically negative, making consensus a globally attracting state~\cite{jadbabaie2003}. Breaking this attraction requires either heterogeneous coupling, noise, or structural frustration~\cite{odor2024}.

We propose a mechanism based on \textbf{doubt-modulated coupling} in extended FitzHugh--Nagumo (FHN) networks. Each node carries a third state variable $u \in [0, 1]$ (the \textit{doubt} variable) that continuously modulates the sign and magnitude of inter-node coupling. When $u > 0.5$, coupling becomes repulsive, implementing frustrated synchronization~\cite{odor2024} at the single-node level. Combined with a fraction $\eta$ of nodes receiving inverted external stimulus (\textit{heretic nodes}), this architecture prevents consensus collapse from homogeneous initial conditions without requiring noise as the sole symmetry-breaking mechanism (heretics provide structural asymmetry; noise enables the Cold Start Protocol described in Section~2.3).

Our contributions are: (1)~a formal specification of the doubt-modulated FHN model with five anti-synchronization mechanisms (doubt-modulated coupling, heretic nodes, $1/\sqrt{N}$ scaling, adaptive doubt speed, and degree-normalized coupling); (2)~empirical demonstration of sustained attractor diversity ($H_{\text{stable}} \approx 0.92$ bits) on regular lattices; (3)~a systematic investigation of degree-normalized coupling on scale-free networks, revealing an abrupt phase transition at algebraic connectivity $\lambda_2 \approx 2$--$3$ beyond which no per-node normalization preserves diversity; and (4)~identification of the optimal power-law exponent $\gamma^*(m)$ for $D_{\text{eff}}(i) = D/\deg(i)^\gamma$ as a function of network density.

% ============================================================================
% 2. MODEL
% ============================================================================
\section{Model}

\subsection{FitzHugh--Nagumo Dynamics with Doubt}

Each node $i$ in a network of $N$ units evolves according to a modified FitzHugh--Nagumo system with a third variable $u_i$ (doubt):
\begin{align}
\frac{dv_i}{dt} &= v_i - \frac{v_i^3}{5} - w_i + I_{\text{ext},i} - \alpha_{\text{self}} \tanh(v_i) + \eta_v(t), \label{eq:dv} \\
\frac{dw_i}{dt} &= \varepsilon(v_i + a - bw_i) + dw_{\text{plasticity},i}, \label{eq:dw} \\
\frac{du_i}{dt} &= \frac{\varepsilon_u}{\tau_u}\bigl(k_u \sigma_{\text{local},i} + \sigma_{\text{baseline}} - u_i\bigr), \label{eq:du}
\end{align}
where $v_i$ is the activation potential, $w_i$ is the recovery variable, $u_i \in [0, 1]$ is the doubt variable (clamped), $\eta_v(t) \sim \mathcal{N}(0, \sigma_{\text{noise}}^2)$ is Gaussian noise, $\alpha_{\text{self}}$ is the self-inhibition gain, and $\bar{v}_{\mathcal{N}(i)} = |\mathcal{N}(i)|^{-1} \sum_{j \in \mathcal{N}(i)} v_j$ is the mean activation of neighbors of node~$i$. The local disagreement $\sigma_{\text{local},i} = |\bar{v}_{\mathcal{N}(i)} - v_i|$ measures the mismatch between a node's state and its neighborhood mean. The plasticity term $dw_{\text{plasticity},i} = \varepsilon_{\text{plast}} \cdot \sigma_{\text{local},i} \cdot (w_{\text{target}} - w_i)$ provides activity-dependent inhibition (a novel element of this model, not present in standard FHN formulations).

The doubt dynamics admit an adaptive extension where the effective learning rate $\varepsilon_{u,\text{eff}}(i)$ replaces the fixed $\varepsilon_u$ (Section~\ref{sec:adaptive}).

The external input for each node depends on its type:
\begin{equation}
I_{\text{ext},i} = \begin{cases}
+I_{\text{stimulus}} + I_{\text{coupling},i} & \text{if node } i \text{ is normal,} \\
-I_{\text{stimulus}} + I_{\text{coupling},i} & \text{if node } i \text{ is a heretic node.}
\end{cases}
\label{eq:stimulus}
\end{equation}

\subsection{Doubt-Modulated Coupling}

The coupling term implements frustrated synchronization via a shifted-tanh kernel:
\begin{equation}
I_{\text{coupling},i} = \frac{D}{\sqrt{N}} \cdot w_i(u_i) \cdot \bigl(\bar{v}_{\mathcal{N}(i)} - v_i\bigr),
\label{eq:coupling}
\end{equation}
where the coupling weight is
\begin{equation}
w_i(u_i) = \tanh\bigl(\pi(0.5 - u_i)\bigr) + \delta, \qquad \delta = 0.01.
\label{eq:sigmoid}
\end{equation}
The factor $\pi$ normalizes the slope at the transition point $u = 0.5$ to $-\pi \approx -3.14$, ensuring a sharp polarity flip. For $u_i < 0.5$, $w_i > 0$ (attractive coupling: the node tends toward its neighbors). For $u_i > 0.5$, $w_i < 0$ (repulsive coupling: the node actively diverges from its neighbors). The offset $\delta$ ensures $w_i(0.5) = \delta > 0$, eliminating the dead zone at the transition point. This kernel can be interpreted as a continuous analogue of frustrated interactions in spin glass systems~\cite{toulouse1977,odor2024}.

\subsection{Heretic Nodes and Structural Heterogeneity}

A fraction $\eta$ of nodes are designated as \textit{heretic nodes}: they receive their external stimulus with inverted polarity (Eq.~\ref{eq:stimulus}). This creates a permanent structural asymmetry analogous to random-field disorder in spin systems~\cite{nattermann1998}. On a $k$-regular lattice with heretic fraction $\eta$, the probability that a normal node is adjacent to at least one heretic is:
\begin{equation}
P_{\text{contact}} = 1 - (1-\eta)^k.
\label{eq:percolation}
\end{equation}
For $k=4$ (2D grid) and $\eta = 0.15$: $P_{\text{contact}} \approx 0.48$, meaning nearly half of all normal nodes experience direct heretic influence. This is consistent with the empirically observed critical threshold on regular lattices (Section~\ref{sec:results_lattice}).

\textbf{Cold Start Protocol.} All simulations use homogeneous initial conditions ($v_i = w_i = 0$, $H(0) = 0$). Diversity must \emph{emerge} from noise and heretic-driven symmetry breaking, not from favorable initialization.

\subsection{Adaptive Extensions}
\label{sec:adaptive}

\textbf{Adaptive doubt speed.} The fixed doubt learning rate $\varepsilon_u$ in Eq.~\ref{eq:du} can be replaced by a per-node adaptive rate modulated by local disagreement:
\begin{equation}
\varepsilon_{u,\text{eff}}(i) = \varepsilon_u \cdot \min\!\bigl(1 + \alpha_s \cdot \sigma_{\text{local}}(i),\; C_{\text{cap}}\bigr),
\label{eq:meta_doubt}
\end{equation}
with surprise gain $\alpha_s = 2.0$ and cap $C_{\text{cap}} = 5.0$. When a node's neighborhood contradicts its internal state ($\sigma_{\text{local}}$ large), its doubt updates faster. Setting $\alpha_s = 0$ recovers the fixed-rate dynamics of Eq.~\ref{eq:du}.

\textbf{Topological rewiring.} On networks with explicit adjacency matrices, nodes with $u_i > u_{\text{rewire}}$ may disconnect from their most consensual neighbor and reconnect to a random non-neighbor (degree-preserving, cooldown-regulated). This was found insufficient to resolve diversity collapse on scale-free networks alone (Section~\ref{sec:topo_regimes}).

\textbf{Degree-normalized coupling.} To address hub-induced synchronization on heterogeneous networks, the uniform coupling $D/\sqrt{N}$ is replaced with per-node weights:
\begin{equation}
D_{\text{eff}}(i) = \frac{D}{\deg(i)^\gamma}, \qquad \gamma \in [0, 1],
\label{eq:degree_norm}
\end{equation}
normalized so that $\langle D_{\text{eff}} \rangle = D/\sqrt{N}$.

\subsection{Simulation Protocol}

All simulations use Euler integration with $\Delta t = 0.05$, validated for numerical stability to $20{,}000+$ steps via systematic $\Delta t$ sweep ($\Delta t \in [0.01, 0.10]$). The sustained attractor entropy $H_{\text{stable}}$ is computed as the mean Shannon entropy over the last 25\% of each run. Entropy is calculated over 5 discrete states defined by $v$-thresholds (Table~\ref{tab:states}).

\begin{table}[ht]
\centering
\caption{Cognitive state classification by activation potential $v$.}
\label{tab:states}
\begin{tabular}{lcc}
\toprule
\textbf{State} & \textbf{Range} & \textbf{Interpretation} \\
\midrule
1 (Strong negative) & $v < -1.2$ & Strong opposition \\
2 (Weak negative)   & $-1.2 \leq v < -0.4$ & Mild opposition \\
3 (Neutral)         & $-0.4 \leq v \leq 0.4$ & Uncertain / Deliberative \\
4 (Weak positive)   & $0.4 < v \leq 1.2$ & Mild agreement \\
5 (Strong positive) & $v > 1.2$ & Strong agreement \\
\bottomrule
\end{tabular}
\end{table}

Reference parameters: $a = 0.7$, $b = 0.8$, $\varepsilon = 0.08$, $D = 0.15$, $\alpha_{\text{self}} = 0.15$, $\sigma_{\text{noise}} = 0.05$, $\sigma_{\text{baseline}} = 0.05$, $\eta = 0.15$ (see Appendix~\ref{app:parameters} for all parameters). Topologies tested: 2D periodic lattice (stencil Laplacian) and Barab\'asi--Albert ($m \in \{1, \ldots, 15\}$, $N = 100$). Watts-Strogatz and Erd\H{o}s--R\'enyi topologies were used during development for qualitative robustness checks but are not systematically reported here. Seeds and repetitions are reported per experiment.

% ============================================================================
% 3. THEORETICAL ANALYSIS
% ============================================================================
\section{Theoretical Analysis}

\subsection{Existence of Limit Cycles}

For a single decoupled unit ($D = 0$) with $u$ treated as a quasi-static parameter ($\varepsilon_u \ll 1$), the $(v, w)$ subsystem reduces to a modified FHN oscillator. By the Poincar\'e--Bendixson theorem:
\begin{enumerate}
\item The self-inhibition term $-\alpha_{\text{self}} \tanh(v)$ ensures boundedness of $|v|$, and $b > 0$ ensures stability of $w$. A compact positively invariant set $\Omega$ exists.
\item The unique equilibrium is unstable for the standard parameter range (verified numerically).
\item By Poincar\'e--Bendixson, the absence of a stable equilibrium in $\Omega$ guarantees the existence of a limit cycle.
\end{enumerate}

For the coupled system, the Floquet multipliers of the linearized dynamics around the synchronized manifold satisfy $|\mu_k| \gg 1$ for all transverse modes $k > 0$ under numerical evaluation (Section~\ref{sec:comp_analysis}), indicating that the synchronized state is unstable.

\subsection{Heretic Percolation Threshold}

The heretic fraction $\eta = 0.15$ is an \textbf{empirical observation}, not a derived constant. Its effectiveness on regular lattices admits a qualitative interpretation via Eq.~\ref{eq:percolation}: at $\eta = 0.15$ on a 4-regular graph, 48\% of normal nodes are adjacent to at least one heretic, providing sufficient local frustration to prevent consensus nucleation.

An analogy with the random-field Ising model~\cite{nattermann1998} suggests that the critical heretic concentration should scale inversely with mean degree: $\eta_c \approx C / \langle k \rangle$. Fitting $C$ to match the observed threshold on 2D lattices ($\langle k \rangle = 4$, $\eta_c \approx 0.15$) gives $C \approx 0.6$. This is an empirical fit, not a derivation, and should be treated as a guideline. In particular, on scale-free networks, the degree heterogeneity invalidates the mean-field assumption underlying this estimate.

\subsection{Entropy Lower Bound}

We sketch a proof that $H > 0$ whenever $\eta > 0$. Suppose the system reaches consensus ($v_i = v^*$ for all $i$). Then the Laplacian coupling vanishes. Normal nodes receive $+I_{\text{stimulus}}$ while heretics receive $-I_{\text{stimulus}}$, driving them toward different equilibria. This contradicts the consensus assumption, so consensus is not an attractor when $\eta > 0$. A quantitative lower bound via mean-field analysis yields:
\begin{equation}
H \geq -\eta \log_2 \eta - (1 - \eta) \log_2(1 - \eta) \approx 0.61 \text{ bits for } \eta = 0.15,
\label{eq:entropy_bound}
\end{equation}
representing the mixing entropy of the two subpopulations. The observed $H_{\text{stable}} \approx 0.92$ exceeds this bound, consistent with intra-group diversity contributing additional entropy.

% ============================================================================
% 4. RESULTS
% ============================================================================
\section{Results}
\label{sec:results}

\subsection{Diversity on Regular Lattices}
\label{sec:results_lattice}

Table~\ref{tab:scaling} reports the sustained attractor entropy across network scales on 2D periodic lattices. All runs use the Cold Start Protocol ($v = w = 0$), $I_{\text{stimulus}} = 0$, $\eta = 0.15$, 1000 steps.

\begin{table}[ht]
\centering
\caption{Scaling performance on 2D periodic lattices (Cold Start, $\eta = 0.15$, 7 seeds per configuration).}
\label{tab:scaling}
\begin{tabular}{lccc}
\toprule
\textbf{Metric} & \textbf{$4 \times 4$} & \textbf{$10 \times 10$} & \textbf{$25 \times 25$} \\
\midrule
$N$                   & 16   & 100   & 625   \\
$H_{\text{stable}}$   & $0.95 \pm 0.06$ & $0.92 \pm 0.04$  & $0.89 \pm 0.03$  \\
Distinct states        & 3    & 4     & 4     \\
Mean doubt $\bar{u}$   & 0.051 & 0.049 & 0.052 \\
\bottomrule
\end{tabular}
\end{table}

The system sustains $H_{\text{stable}} > 0.85$ across all tested scales. Note: the initial 200--500 time steps exhibit transient entropy peaks ($H_{\text{peak}} \approx 1.5$) before converging to the sustained attractor. All values in this paper report $H_{\text{stable}}$ unless explicitly noted.

\subsection{Component Ablation}

Table~\ref{tab:ablations} quantifies the contribution of each mechanism to diversity preservation.

\begin{table}[ht]
\centering
\caption{Ablation study ($10 \times 10$ lattice, Cold Start, 3000 steps, 5 seeds). $H_{\text{stable}}$ is the mean over the last 25\% of each run.}
\label{tab:ablations}
\begin{tabular}{lcc}
\toprule
\textbf{Configuration} & $H_{\text{stable}}$ & \textbf{Collapse?} \\
\midrule
Full model & $0.92 \pm 0.04$ & No \\
No heretics ($\eta = 0$) & $< 0.05$ & Yes \\
No doubt ($u \equiv 0$) & $0.41 \pm 0.12$ & Partial \\
No noise ($\sigma_{\text{noise}} = 0$) & $0.00$ & Yes (Cold Start) \\
No coupling ($D = 0$) & $0.38 \pm 0.08$ & No (isolated oscillators) \\
\bottomrule
\end{tabular}
\end{table}

The heretic mechanism is \textbf{necessary}: removing it causes consensus collapse under Cold Start. The doubt variable is \textbf{enhancing}: without it, coupling is always attractive, and diversity is sustained only by noise-driven fluctuations ($H_{\text{stable}} \approx 0.41$).

\subsection{Topological Normalization Regimes}
\label{sec:topo_regimes}

We systematically sweep the Barab\'asi--Albert attachment parameter $m \in \{1, 2, 3, 4, 5, 6, 7, 8, 10, 15\}$ with $N = 100$, 5 seeds per configuration, 3000 steps, and $I_{\text{stimulus}} = 0$. Two coupling modes are compared: uniform ($D/\sqrt{N}$, i.e.\ $\gamma = 0$) and degree-linear ($D/\deg(i)$, i.e.\ $\gamma = 1$). Note: Table~\ref{tab:alpha_sweep} extends this analysis to intermediate $\gamma$ values.

\begin{table}[ht]
\centering
\caption{Sustained entropy $H_{\text{stable}}$ by BA attachment parameter $m$ and coupling normalization. ``$\star$'' denotes $H > 0.5$.}
\label{tab:ba_m_sweep}
\begin{tabular}{rrrrl}
\toprule
$m$ & $H_{\text{uniform}}$ & $H_{\text{deg.\ linear}}$ & Spectral gap $\lambda_2$ & Effective normalization \\
\midrule
1  & $0.86\;\star$ & 0.00 & 0.03 & uniform \\
2  & 0.03 & $0.69\;\star$ & 0.63 & degree-linear \\
3  & 0.00 & $0.83\;\star$ & 1.41 & degree-linear \\
4  & 0.00 & 0.23 & 2.21 & degree-linear (marginal) \\
5  & 0.02 & 0.00 & 2.99 & neither \\
6  & 0.04 & 0.00 & 3.92 & neither \\
7  & 0.08 & 0.00 & 5.05 & neither \\
8  & 0.15 & 0.00 & 5.86 & uniform (marginal) \\
10 & 0.21 & 0.00 & 7.70 & uniform (weak) \\
15 & 0.37 & 0.00 & 12.55 & uniform (weak) \\
\bottomrule
\end{tabular}
\end{table}

Two \textbf{regime transitions} are observed:
\begin{enumerate}
\item $m = 1 \to m = 2$: crossover from uniform-dominant (tree topology) to degree-linear-dominant (sparse scale-free).
\item $m = 4 \to m = 5$: abrupt transition to a dead zone where \emph{neither normalization sustains diversity}.
\end{enumerate}

The spectral gap $\lambda_2$ (algebraic connectivity of the graph Laplacian) is the strongest predictor: degree-linear normalization succeeds only when $\lambda_2 < 2$ (sparse, poorly connected networks).

\subsection{Power-Law Exponent Sweep}
\label{sec:alpha_sweep}

To test whether an intermediate exponent $\gamma$ in $D_{\text{eff}}(i) = D/\deg(i)^\gamma$ (Eq.~\ref{eq:degree_norm}) can bridge the dead zone identified in Table~\ref{tab:ba_m_sweep}, we sweep $\gamma \in [0, 1]$ across $m \in \{2, 3, 4, 5, 6, 8, 10\}$ (3 seeds, 3000 steps, $I_{\text{stimulus}} = 0$).

\begin{table}[ht]
\centering
\caption{Optimal exponent $\gamma^*$ and maximum achievable entropy by BA attachment parameter. Table~\ref{tab:ba_m_sweep} tests only $\gamma \in \{0, 1\}$; this table sweeps $\gamma$ over $[0, 1]$ in steps of 0.1, revealing that the optimal $\gamma^* < 1$ for $m = 2$.}
\label{tab:alpha_sweep}
\begin{tabular}{rrrl}
\toprule
$m$ & $\gamma^*$ & $\max H_{\text{stable}}$ & Assessment \\
\midrule
2  & 0.7 & 0.99 & Near-optimal diversity \\
3  & 0.9 & 0.89 & Strong diversity \\
4  & 1.0 & 0.37 & Below threshold \\
$\geq 5$ & --- & $< 0.03$ & Complete failure \\
\bottomrule
\end{tabular}
\end{table}

The transition is \textbf{abrupt}: $H_{\text{stable}}$ drops from 0.89 ($m = 3$) to $< 0.03$ ($m = 5$) with no intermediate regime. The optimal $\gamma^*$ decreases with increasing $m$, consistent with the intuition that sparser networks require stronger per-node correction. For $m = 2$, degree-linear ($\gamma = 1.0$) actually \emph{over-corrects}: the optimum at $\gamma = 0.7$ achieves $H = 0.99$, exceeding even lattice performance.

\textbf{Negative result.} For $m \geq 5$, no $\gamma \in [0, 1]$ achieves $H_{\text{stable}} > 0.03$. Degree-based normalization is fundamentally insufficient when graph connectivity (measured by $\lambda_2$) exceeds $\sim 2$--$3$, because coupling signals bypass locally-corrected hubs through redundant paths.

\subsection{Comparative Benchmarks}

Table~\ref{tab:benchmarks} compares the model against classical collective dynamics models on $10 \times 10$ lattices, 1000 steps, biased stimulus phase ($I_{\text{stim}} = 1.0$ for steps 200--800).

\begin{table}[ht]
\centering
\caption{Benchmark comparison ($N = 100$, $10 \times 10$ lattice). $H_{\text{stable}}$ measured after stimulus removal.}
\label{tab:benchmarks}
\begin{tabular}{lcc}
\toprule
\textbf{Model} & $H_{\text{stable}}$ & \textbf{Diversity sustained?} \\
\midrule
Doubt-modulated FHN (this work) & $0.92 \pm 0.04$ & Yes \\
Kuramoto ($K = 2$) & $< 0.01$ & No (phase-locked) \\
Voter model & $0.00$ & No (absorbing state) \\
Distributed consensus & $0.00$ & No (by design) \\
\bottomrule
\end{tabular}
\end{table}

% ============================================================================
% 5. COMPUTATIONAL ANALYSIS
% ============================================================================
\section{Computational Analysis}
\label{sec:comp_analysis}

\subsection{Complexity}

On stencil-based lattices, each time step requires $\approx 56N$ floating-point operations (5-point Laplacian: $9N$; sigmoid: $4N$; FHN dynamics and doubt: $43N$), yielding $O(N)$ per-step complexity. Empirical benchmarks confirm:
\begin{equation}
t_{\text{step}}(\text{ms}) \approx 2.87 \times 10^{-4} N + 0.212,
\label{eq:scaling}
\end{equation}
enabling real-time simulation up to $N \approx 10{,}000$ on commodity hardware. For explicit adjacency matrices, Laplacian computation is $O(N^2)$, but incremental updates reduce per-rewire cost to $O(1)$. For $N > 1000$, the adjacency matrix is automatically converted to compressed sparse row (CSR) format, yielding $455\times$ memory reduction at $N = 5000$.

\subsection{Numerical Stability}

The model uses first-order Euler integration with $\Delta t = 0.05$. A systematic sweep ($\Delta t \in \{0.01, 0.02, 0.05, 0.10\}$, 20{,}000 steps) confirmed that $\Delta t \leq 0.05$ is numerically stable: after an initial transient convergence ($\sim$500 steps), entropy standard deviation drops to $\sim$0.016 and remains flat. Higher-order methods (Runge--Kutta 4) improve precision but are not required for stability.

\subsection{Floquet Analysis and Lyapunov Exponents}

Under periodic forcing $I(t) = A\sin(\omega t)$ ($A = 0.3$, $\omega = 0.5$), we computed the Floquet multipliers of the linearized dynamics around the synchronized manifold. For the synchronized state to be stable, transverse Floquet multipliers must satisfy $|\mu| < 1$. Here, all transverse multipliers satisfy $|\mu| \gg 1$ (magnitudes $\sim 10^3$--$10^4$), confirming instability of the synchronized state and resistance to periodic entrainment. The maximum Lyapunov exponent $\lambda_{\max} \approx 0$ indicates quasi-periodic dynamics---neither chaotic ($\lambda_{\max} > 0$) nor convergent ($\lambda_{\max} < 0$).

% ============================================================================
% 6. DISCUSSION
% ============================================================================
\section{Discussion}

\subsection{Frustrated Synchronization}

The doubt-modulated coupling mechanism implements frustrated synchronization at the single-node level: when a node's doubt exceeds $u = 0.5$, its coupling flips from attractive to repulsive. This is analogous to frustrated spin systems, where competing interactions prevent ground-state ordering~\cite{toulouse1977,odor2024}. The heretic fraction provides an additional source of structural frustration through permanent stimulus asymmetry, analogous to quenched disorder in random-field models~\cite{nattermann1998}.

\subsection{The Topological Boundary}

The most significant finding is the abrupt phase transition at $m \approx 5$ (or equivalently, $\lambda_2 \approx 2$--$3$). This result has a clear physical interpretation: degree-based normalization corrects coupling strength at \emph{individual} hubs, but when path redundancy is high ($\lambda_2 \gg 1$), the synchronization signal reaches each node through multiple independent paths, bypassing the corrected hub. The problem is not the strength of any single connection but the multiplicity of alternative routes.

This suggests that resolving diversity collapse on dense scale-free networks requires normalization schemes that account for \textbf{global graph structure}---such as eigenvector centrality-based coupling ($D_{\text{eff}}(i) \propto 1/c_i^{\text{eigen}}$) or symmetric Laplacian whitening. These approaches remain open for future work.

\subsection{Scope and Applicability}

The doubt-modulated FHN model was originally motivated by deliberative scenarios in which premature consensus is undesirable (e.g., citizens' assemblies, collective decision-making). While the model captures the dynamical requirements for diversity preservation, it has not been validated against real-world deliberative data. The results presented here are purely computational and should be understood as properties of the dynamical system, not as claims about social behavior.

% ============================================================================
% 7. CONCLUSION
% ============================================================================
\section{Conclusion}

We have presented a FitzHugh--Nagumo network model in which doubt-modulated coupling and structural heretic nodes sustain attractor diversity from homogeneous initial conditions. On regular lattices, the model reliably achieves $H_{\text{stable}} \approx 0.92$ bits with $\eta = 0.15$ heretic fraction. Degree-normalized coupling ($D/\deg(i)^\gamma$) extends this to sparse scale-free networks ($m \leq 3$), with optimal exponent $\gamma^*(m)$ varying from 0.7 ($m = 2$) to 0.9 ($m = 3$).

The principal negative result is that for Barab\'asi--Albert networks with $m \geq 5$, no degree-based normalization preserves diversity. This phase transition, correlated with the algebraic connectivity $\lambda_2$ exceeding $\sim$2--3, establishes a fundamental boundary for local coupling corrections and motivates future investigation of spectrally-informed normalization.

% ============================================================================
% 8. LIMITATIONS & TRANSPARENCY
% ============================================================================
\section{Limitations and Transparency}

\subsection{Scientific Limitations}
\begin{itemize}
\item \textbf{Numerical integration:} All results use first-order Euler integration. While stable at $\Delta t = 0.05$, quantitative values may shift under higher-order schemes.
\item \textbf{Heretic threshold:} The $\eta = 0.15$ value is empirically effective on 2D lattices ($k = 4$) but has not been validated as optimal across arbitrary topologies.
\item \textbf{State discretization:} The 5-state classification (Table~\ref{tab:states}) is a design choice. Entropy values depend on bin boundaries and should not be compared across different discretization schemes.
\item \textbf{Scale:} The model has been tested up to $N = 2500$. Behavior at $N \gg 10{,}000$ is not validated.
\item \textbf{Topological phase transition:} The transition at $m \approx 5$ (or $\lambda_2 \approx 2$--$3$) is empirically observed on $N = 100$ BA networks. Finite-size effects may shift this boundary at larger $N$.
\end{itemize}

\subsection{Methodological Transparency}
This work was developed through a multi-agent collaborative research methodology involving multiple AI systems (Claude, GPT, Gemini) under continuous human guidance. All scientific claims, parameter choices, and publication decisions were made by the human author. AI contributions were advisory and generative. Complete session transcripts and intermediate versions are maintained in the public repository. The model underwent systematic adversarial testing via automated analysis of failure modes (December 2025--April 2026), leading to the corrections of entropy measurement (transient vs.\ sustained values) and the development of degree-normalized coupling.

% ============================================================================
% REFERENCES
% ============================================================================
\begin{thebibliography}{99}

% Consensus & Synchronization
\bibitem{olfati2007} Olfati-Saber, R., Fax, J.A., Murray, R.M., \emph{Consensus and cooperation in networked multi-agent systems}, Proceedings of the IEEE, 95(1):215-233, 2007.
\bibitem{strogatz2000} Strogatz, S.H., \emph{From Kuramoto to Crawford: Exploring the onset of synchronization in populations of coupled oscillators}, Physica D, 143(1-4):1-20, 2000.
\bibitem{acebrn2005} Acebr\'on, J.A., et al., \emph{The Kuramoto model: A simple paradigm for synchronization phenomena}, Reviews of Modern Physics, 77(1):137, 2005.

% FitzHugh-Nagumo
\bibitem{fitzhugh1961} FitzHugh, R., \emph{Impulses and physiological states in theoretical models of nerve membrane}, Biophysical Journal, 1(6):445-466, 1961.
\bibitem{nagumo1962} Nagumo, J., Arimoto, S., Yoshizawa, S., \emph{An active pulse transmission line simulating nerve axon}, Proceedings of the IRE, 50(10):2061-2070, 1962.
\bibitem{izhikevich2007} Izhikevich, E.M., \emph{Dynamical Systems in Neuroscience: The Geometry of Excitability and Bursting}, MIT Press, 2007.

% Network Science
\bibitem{barabasi1999} Barab\'asi, A.-L., Albert, R., \emph{Emergence of scaling in random networks}, Science, 286(5439):509-512, 1999.
\bibitem{watts1998} Watts, D.J., Strogatz, S.H., \emph{Collective dynamics of `small-world' networks}, Nature, 393(6684):440-442, 1998.

% Opinion Dynamics
\bibitem{castellano2009} Castellano, C., Fortunato, S., Loreto, V., \emph{Statistical physics of social dynamics}, Reviews of Modern Physics, 81(2):591, 2009.
\bibitem{hegselmann2002} Hegselmann, R., Krause, U., \emph{Opinion dynamics and bounded confidence models, analysis, and simulation}, Journal of Artificial Societies and Social Simulation, 5(3), 2002.

% Multi-Agent Systems
\bibitem{jadbabaie2003} Jadbabaie, A., Lin, J., Morse, A.S., \emph{Coordination of groups of mobile autonomous agents using nearest neighbor rules}, IEEE Transactions on Automatic Control, 48(6):988-1001, 2003.

% Nonlinear Dynamics
\bibitem{pikovsky2001} Pikovsky, A., Rosenblum, M., Kurths, J., \emph{Synchronization: A Universal Concept in Nonlinear Sciences}, Cambridge University Press, 2001.

% Frustrated Synchronization & Spin Glasses
\bibitem{toulouse1977} Toulouse, G., \emph{Theory of the frustration effect in spin glasses: I}, Communications on Physics, 2:115--119, 1977.
\bibitem{odor2024} \'Odor, G., Deng, S., Kelling, J., \emph{Frustrated Synchronization of the Kuramoto Model on Complex Networks}, Entropy, 26(12):1074, 2024.

% Random-Field Ising Model
\bibitem{nattermann1998} Nattermann, T., \emph{Theory of the Random Field Ising Model}, in \emph{Spin Glasses and Random Fields}, ed.\ A.P.\ Young, World Scientific, pp.\ 277--298, 1998.

% Adaptive Networks
\bibitem{gross2008} Gross, T., Blasius, B., \emph{Adaptive coevolutionary networks: a review}, Journal of the Royal Society Interface, 5(20):259-271, 2008.

% AI Methodology
\bibitem{bommasani2021} Bommasani, R., et al., \emph{On the Opportunities and Risks of Foundation Models}, arXiv:2108.07258, 2021.
\bibitem{park2023} Park, J.S., et al., \emph{Generative Agents: Interactive Simulacra of Human Behavior}, Proceedings of the 36th Annual ACM Symposium on UIST, 2023.

\end{thebibliography}

% ============================================================================
% APPENDICES
% ============================================================================
\appendix

\section{Full Algorithm}
\label{app:algorithm}

\begin{algorithm}[ht]
\caption{Time-step update for the doubt-modulated FHN network.}
\begin{algorithmic}[1]
\FOR{each time-step $t$}
    \STATE \COMMENT{Optional: Doubt-driven rewiring (explicit graphs only)}
    \FOR{each unit $i$ with $u_i > u_{\text{rewire}}$ and cooldown expired}
        \STATE $j^* \leftarrow \arg\min_{k \in \mathcal{N}(i)} |v_i - v_k|$
        \STATE $k \leftarrow \text{RandomNonNeighbor}(i)$
        \STATE Disconnect $i \leftrightarrow j^*$; Connect $i \leftrightarrow k$
    \ENDFOR
    \FOR{each unit $i$}
        \STATE $\sigma_{\text{local},i} \leftarrow |\bar{v}_{\mathcal{N}(i)} - v_i|$
        \STATE $\eta_v \leftarrow \mathcal{N}(0, \sigma_{\text{noise}}^2)$
        \STATE $w_i^{\text{coup}} \leftarrow \tanh(\pi(0.5 - u_i)) + \delta$
        \STATE $I_{\text{coup}} \leftarrow D_{\text{eff}}(i) \cdot w_i^{\text{coup}} \cdot (\bar{v}_{\mathcal{N}(i)} - v_i)$
        \STATE $I_{\text{ext}} \leftarrow s_i \cdot I_{\text{stimulus}} + I_{\text{coup}}$ \COMMENT{$s_i = -1$ for heretics}
        \STATE $v_i \leftarrow v_i + (v_i - v_i^3/5 - w_i + I_{\text{ext}} - \alpha_{\text{self}} \tanh(v_i) + \eta_v) \cdot \Delta t$
        \STATE $w_i \leftarrow w_i + (\varepsilon(v_i + a - bw_i) + dw_{\text{plast},i}) \cdot \Delta t$
        \STATE $\varepsilon_{u,\text{eff}} \leftarrow \varepsilon_u \cdot \min(1 + \alpha_s \sigma_{\text{local},i},\; C_{\text{cap}})$
        \STATE $u_i \leftarrow \text{clamp}\!\bigl(u_i + \frac{\varepsilon_{u,\text{eff}}}{\tau_u}(k_u \sigma_{\text{local},i} + \sigma_{\text{base}} - u_i) \cdot \Delta t,\; 0,\; 1\bigr)$
    \ENDFOR
    \STATE Compute $H(t)$ from state distribution (Table~\ref{tab:states})
\ENDFOR
\end{algorithmic}
\end{algorithm}

\section{Reference Parameters}
\label{app:parameters}

\begin{table}[ht]
\centering
\caption{Default parameter values. All parameters are configurable via YAML.}
\begin{tabular}{llrl}
\toprule
\textbf{Group} & \textbf{Parameter} & \textbf{Value} & \textbf{Description} \\
\midrule
Dynamics & $a$ & 0.7 & FHN offset \\
         & $b$ & 0.8 & Recovery rate \\
         & $\varepsilon$ & 0.08 & Time-scale separation \\
         & $\alpha_{\text{self}}$ & 0.15 & Self-inhibition gain \\
         & $\Delta t$ & 0.05 & Integration step \\
\midrule
Coupling & $D$ & 0.15 & Base coupling strength \\
         & $\delta$ & 0.01 & Sigmoid offset \\
         & $\eta$ & 0.15 & Heretic fraction \\
\midrule
Doubt    & $\varepsilon_u$ & 0.02 & Doubt learning rate \\
         & $\tau_u$ & 10.0 & Doubt time constant \\
         & $k_u$ & 1.0 & Local disagreement gain \\
         & $\sigma_{\text{baseline}}$ & 0.05 & Minimum doubt \\
         & $u_{\text{clamp}}$ & $[0, 1]$ & Doubt range \\
\midrule
Meta-doubt & $\alpha_s$ & 2.0 & Surprise gain \\
           & $C_{\text{cap}}$ & 5.0 & Acceleration cap \\
\midrule
Noise    & $\sigma_{\text{noise}}$ & 0.05 & State noise std \\
\midrule
Rewiring & $u_{\text{rewire}}$ & 0.8 & Rewire threshold \\
         & Cooldown & 50 & Steps between rewires \\
\bottomrule
\end{tabular}
\end{table}

\end{document}
